% Supplemental figures for Document S1 (Gate C): files, titles (the build docstrings), the main figure each relates to, and
% legends drawn from the builds' panel lists and source data. Not counted toward the 7,000 words.
\setfigfile{S1}{../figures/v4/supplement/s1.pdf}\relatedto{S1}{1}
\setfigfile{S2}{../figures/v4/supplement/s2.pdf}\relatedto{S2}{2}
\setfigfile{S3}{../figures/v4/supplement/s3.pdf}\relatedto{S3}{3}
\setfigfile{S4}{../figures/v4/supplement/s4.pdf}\relatedto{S4}{4}
\setfigfile{S5}{../figures/v4/supplement/s5.pdf}\relatedto{S5}{5}
\setfigfile{S6}{../figures/v4/supplement/s6.pdf}\relatedto{S6}{6}
\setfigfile{S7}{../figures/v4/supplement/s7.pdf}\relatedto{S7}{7}

\deflegend{S1}{The model, what training changes, and its validation}{%
(A) Where training acts in the planning model: each cell by the mean absolute change of its cell-type group's output gain, input
gain and threshold over the \val{L09} maps (log color), and the cells of the most consistently changed groups. Scale bar,
100~μm.
(B) Total parameter change against training length, one dot per map; Spearman ρ with a BCa bootstrap interval, permutation
test, Holm-adjusted.
(C) Mean absolute change by cell group and parameter.
(D) Every map's learning curve: solution depth reached against training updates, by difficulty tier.
(E) Updates to solve each map (rows by optimal moves) in the original set and five replicate sets.
(F) Updates per stage for stages entered with or without a deadlocking move available; bars, medians.
(G) The hardest map: training updates of the analyzed lineage in six sets, and from a neutral start with the real wiring, with
its weights shuffled and rewired with degrees kept (one run each); cross, stopped without progress.
(H) The model's parameters at a glance.
(I) Bit-exact replay of every analyzed model on every graphics processor of the three training hosts.}

\deflegend{S2}{The silent trace in every learner: written, held, read}{%
(A) Where the cue writes the trace: the cells that spike differently during the cue (scene 1), and the voltage difference held
at the end of the delay in scenes 1--3, with the cells that carry an adaptation difference. Scale bar, 100~μm.
(B) The cue difference at the decision in five state sets: scene 1 with the readout or the brain and readout trained, scene 2,
and scene 3 in each of its two contexts.
(C) The share of cells above each absolute cue difference, at the end of the delay and at the decision, four learners.
(D) Time constants of the cue difference by state variable, in every learner, timing and context; line, the intrinsic time
constant.
(E) The synaptic share of the decision-time drive on the grid of delay and prefix input, four learners; framed, trained timings.
(F) The decision margin with the transplanted state against the native margin, at every delay, without and with prefix input.
(G) Scene 3 holds its context with the cue: trained histories solved with the own state, the opposite cue's state or the other
context's state; exact tests, Holm- and BH-adjusted.
(H) The cue difference by state variable through the \val{H12}-ms delay, in every learner.}

\deflegend{S3}{No circuit is special, in every scene and under every connection rule}{%
(A) The scene-2 cue ensemble (the \val{D15} of cells with the largest cue difference at the decision) and its internal
connections, beside three of its \val{D14} matched sets drawn by the same rule. Scale bar, 100~μm.
(B) The cell groups of the scene-1 and scene-2 ensembles, which their matched sets share by design, against all cells.
(C) Trained histories solved after silencing the ensemble or cutting its internal synapses, relative to the mean of the matched
sets (gray), by the ensemble's scene and the scene tested.
(D) Removing whole populations: trained histories solved, the network's rate and the giant fiber's looming response, against
healthy.
(E) The signatures of healthy fly circuits used by the disease study (a stable network, antennal-lobe inhibition, looming
escape and the monoamine control) and the task rate of the healthy network under every connection rule.}

\deflegend{S4}{The molecular atlas behind the \gene{SK} signature}{%
(A) Single-gene in silico stains on the unfolded CNS (\gene{SK}, \gene{slo}, \gene{KCNQ}, \gene{Shab}, \gene{Shal} and the
adaptation index): each mapped cell by its atlas unit's z score, with the mapping tier as brightness.
(B) The share of each cell set mapped by cell type, by hemilineage and by hemilineage pair.
(C) Agreement with the connectome, per atlas: the transmitters of the mapped types against the predictions from electron
microscopy, and the map rows present.
(D) Agreement between atlases: the expression correlation of units mapped to the same cell type, against different types.
(E) Marker genes in the mapped units of Kenyon cells, projection neurons and ring neurons, one dot per atlas.
(F) \gene{SK} enrichment of the scene-1 carriers with and without the hemilineage-pair tier, against matched sets and the
mapping null.
(G) The cue's routes to the motor neurons by hemilineage: modules tinted by atlas \gene{SK} (z); gold, types that held the
trace; outlined, motor neurons; width, the cue drive at the first decision; beside it, each hemilineage's \gene{SK} and
carrier share.}

\deflegend{S5}{Access locked to the timing of learning, at every delay}{%
(A) The decision-time cue difference at all seven delays (four prefix windows; scene 1, readout trained), each with its outcome
and its cosine with the trained \val{G06}-ms pattern.
(B) Cues solved at every delay of the dense scan and every prefix input; framed, trained timings.
(C) At untrained timings the push toward the correct target collapses and its direction turns away from the trained push, in
the three readouts of scene 1 in which the collapse survives correction; bars, medians; Holm-adjusted \Pv{}.
(D) The cosine with the \val{G06}-ms pattern within each cell group, by delay.
(E) One trace pattern injected at seven delays, without and with prefix input: the cosine between response directions, and
each response's push; framed, trained delays.
(F) Untrained use after training at fixed or jittered timings, timing by timing (means of \val{L19} replicates per arm), and
their difference.}

\deflegend{S6}{Many routes, fading kicks and a wider window, in detail}{%
(A) The route network at every tenth of removal, random (the first of three seeded draws) and hubs first (cells in descending
order of connections); numbers, cells on routes.
(B) Cells on routes, and the readout cells by shortest route, in every draw and both orders.
(C) Single extra spikes: per kicked cell, the growth of the difference after the spike and the cells still changed at the
decision, at two delays and for both cues, with the cell's direct partners.
(D) Bootstrap resamples of the access window and the cell-loss tolerance, across states over one state; triangles, estimates;
lines, \val{G11} intervals.
(E) Memory under each held-out jitter draw, replicate by replicate; bars, means.
(F) Flow capacity from the visual input to the readout and from the odor input to the mushroom-body output, with Kenyon cells
removed.
(G) Typed routes from the visual input to the readout, intact and with \val{G08} of the other cells removed: the relay types
with the most route traffic, the input and readout types that use them most, and every connection among them above a
threshold; width, walks of two or three synapses through the connection; color and arrowhead, the sender's transmitter.}

\deflegend{S7}{Disease steals access, not memories, in detail}{%
(A) Each cell's change in task rate under the uniform Kv4-loss lesion at half and full dose and under the \gene{Shal} pattern of
the same mean, and the vulnerable cells (the \val{D15} of memory-relevant cells with the largest change) by group.
(B) The healthy circuits match the fly in both brains: sparse Kenyon-cell coding broadened by APL removal, projection-neuron
rates raised by removing the local neurons, and the giant fiber's looming response.
(C) Lesions after learning in every retained case and memory model; equivalence tests (TOST), Holm-adjusted.
(D) Learning under the lesion, run by run; bars, means.
(E) Relearning after the lesion, replicate by replicate, with the readout retrained.
(F) Lifting the lesion late removes access in every run.
(G) Odor memory by delay in every case, read from the mushroom-body output.
(H) The two \gene{Shal} weightings of one mean loss: class means and the global mean.
(I) Above full dose both weightings run away (cells above \val{G03}~Hz); below, spikes per trial against bits of cue per
decision.
(J) Trained histories solved when each superclass's \gene{Shal} weighting is made uniform in turn (descriptive).
(K) The aging operator: its source, the summed potassium-channel expression of old over young flies in neuronal atlas
units\cellcite{davie2018}, and its effect on Kenyon-cell odor responses in two brains.}
